\section{Mental Health and Student Workload}\label{mental-health-and-student-workload}

\stanceinfo{\textbf{CFES Congress 2018} [2018-01-07]}{2024-01-02}

\officialstance{The Canadian Federation of Engineering Students believes that engineering students face negative mental health outcomes as a consequence of the structure and workload demands of their programs. All stakeholders, including faculty, accreditors, and student leaders, have a responsibility to review and amend their practices to give students a healthy and effective learning environment.}

\stanceheading{The Students' Position}
\begin{itemize}
 \item Engineering students are at an elevated risk of negative mental health outcomes, with four themes remaining consistent throughout literature [3];
 \begin{itemize}
  \item That engineering workload is a defining stressor.
  \item Barriers exist that prevent engineering students from seeking help for mental health concerns.
  \item Peers rely more heavily on one another to cope with stress and mental health distress.
  \item The COVID-19 pandemic exacerbated student challenges surrounding mental health.
 \end{itemize}
 \item Engineering students experience high levels of self-reported stress, anxiety and depression with some literature reportings of 28\% moderate to extremely severe stress, 35\% moderate to extremely severe anxiety and 34\% moderate to extremely severe depression from a survey conducted in 2021 [4].
 \item The generalized perception that high stress and poor mental health is normalized and to be expected in engineering is harmful and dismissive of mental health issues for current undergraduate students and students wanting to choose engineering.
 \item Excessive student workload is likely a primary cause of these negative mental health outcomes, and steps should be taken to reduce academic stress related to workload.
 \item Poor mental health is disproportionately prevalent amongst women and other minorities including LGBTQ2S+ people, Indigenous people, and people of colour [4].
 \item Engineering students are unlikely to seek mental health support even if they self-report higher levels of mental health distress due to different factors including physical and cultural factors [5].
\end{itemize}

\stanceheading{Recommendations to Partners, Stakeholders, and Other Entities}
\begin{itemize}
 \item The CFES calls on the Canadian Engineering Accreditation Board (CEAB) to investigate changes to accreditation criteria that better account for the full burden of student workload, and that require access to basic mental health services for all undergraduate engineering students.
 \item The CFES calls on the CEAB to investigate and review the Accreditation Units (AU) with respective hours to ensure they remain feasible within a full student workload.
 \item The CFES calls on the Engineering Deans Canada (EDC) to review the Mental Health Commission of Canada's National Standard for Mental Health and Well-Being for Post-Secondary Standards (MHCC Post Secondary Standards) [6] and implement these standards in their engineering faculties.
 \item The CFES calls on the EDC to hold their instructors accountable for ensuring workload does not harm engineering students in terms of mental health and instead work towards implementing more lifelong learning opportunities, problem based learning opportunities and less high stake academic challenges.
 \item The CFES calls on Engineers Canada and the EDC to partner with the CFES on research related to student mental health and workload, with the intention of clarifying the specific impacts of workload on the mental health of engineering students, and the best procedures for improving student mental health in engineering programs.
 \item The CFES calls on its partnered regional organizations (WESST, ESSCO, QCESO, ACES) to collaborate on future initiatives in student workload and mental health, potentially including a workload tracking system or a second national student survey.
 \item The CFES calls on the membership to evaluate their workload objectively and find gaps/problems within the engineering curriculum at their schools, and elevate the issue to their regional organizations.
\end{itemize}

\stanceheading{What the CFES is doing}
\begin{itemize}
 \item The CFES has conducted a National Student Survey in 2023 to collect data on the mental health of engineering students as well as student workload/perception. The CFES has run sessions and talks on student mental health, with a focus on student initiatives and student volunteer burnout, at its major events including the Canadian Engineering Leadership Conference and the Conference on Diversity in Engineering.
 \item The CFES has created a bank of Mental Health Best Practices, available to all member societies via the external drive.
\end{itemize}

\stanceheading{What the CFES plans to do}
\begin{itemize}
 \item The CFES will continue to facilitate the sharing of best practices related to mental health between its member societies, and update resources for student mental health advocacy when appropriate.
 \item The CFES will pursue further research into the specific consequences of negative student mental health, and best practices for improving mental health at engineering schools.
 \item The CFES will evaluate how it can better promote student mental health internally through the operations and content of its events.
 \item The CFES will follow up with relevant stakeholders including EDC to ensure systematic review is being done regularly of engineering curriculum workload, hours and policies in place to protect students.
 \item The CFES will review and apply MHCC Post Secondary Standards to conference planning, motivation for advocacy, and internal team practices.
\end{itemize}

\newpage
